\chapter{Introduction}

Quadratic integer sequences play a central role in number theory, 
particularly in the study of congruence properties, polygonal numbers, 
and prime number distributions. The present work considers the special 
quadratic sequence
\begin{equation}\label{eq:k_main}
k(n) = \frac{1}{6}(2n^2 + 3n + 1),
\end{equation}
which possesses a remarkable Diophantine property: it assumes 
integer values if and only if $n$ belongs to the congruence classes
\begin{equation}\label{eq:congruence}
n \equiv 1, 5 \pmod{6}.
\end{equation}

These two classes are of fundamental number-theoretic significance, as they
contain all prime numbers greater than~$3$ (Theorem~\ref{satz:primzahlen} in
Chapter~\ref{sec:arithmetik}).

\section*{Motivation and Objectives}

The geometric visualisation of number-theoretic structures has a long 
tradition and has led to significant insights. Well-known examples include:
\begin{itemize}[leftmargin=2em]
\item The \emph{Ulam spiral}~\cite{Stein1964}
\item The \emph{Sacks spiral}~\cite{Sacks2003}
\item Quadratic forms in prime number theory
\end{itemize}

While these classical constructions are primarily empirically motivated, 
the present work pursues an algebraic-geometric approach: 
starting from the quadratic sequence~\eqref{eq:k_main}, a geometric 
structure is \emph{derived}, not merely visualised.

\clearpage

\section*{Main Results}

The central results of this work are:
\begin{enumerate}
\item Alternating minor and major steps with closed-form expressions
\item Circle circumferences with constant radius difference $\Delta R = \frac{12}{\pi}$
\item Archimedean spiral $r(\theta) = \frac{6}{\pi^2}\theta$
\item Parametrisation via $p$: $r(p)=\frac{6}{\pi}\sqrt{p/3}$, $\theta(p)=\pi\sqrt{p/3}$
\item Exact arc length integral
\item Conical helix as 3D embedding
\item All primes $p>3$ appear as a distinguished discrete subset of the helix (corollary of the congruence class structure)
\end{enumerate}

\section*{Structure of the Paper}

The paper is organised into five parts. \textbf{Part~I}
(Chapters~\ref{sec:arithmetik}--\ref{sec:distanz_sequenzen}) develops the
arithmetic foundations: congruence condition, difference structure, and
Euclidean distances. \textbf{Part~II} (Chapter~\ref{sec:spirale}) derives
the Archimedean spiral from the uniqueness of the pairing.
\textbf{Part~III} (Chapters~\ref{sec:bogenlaenge}--\ref{sec:bogenlaenge_parabel})
treats parametrisation, arc length, and Taylor approximation.
\textbf{Part~IV} (Chapters~\ref{sec:kegel}--\ref{ch:primes}) constructs the
conical helix and derives the prime number corollary. \textbf{Part~V}
(Chapters~\ref{sec:visualisierung}--\ref{sec:kegelstrahlen}) is devoted to the
exact $Q$-parametrisation, visualisation, and cone geometry.
Chapter~\ref{sec:zusammenfassung} provides a summary.

\section*{Notation}

\begin{itemize}[leftmargin=2em]
\item $k(n)$: quadratic mean value sequence
\item $\Delta_n$: step sizes
\item $a_n$: generator function
\item $\theta$: angular parameter
\item $p$: function value parameter
\item $r(\theta)$: radius function
\item $\gamma(\theta)$: helix parametrisation
\end{itemize}
